\scanpage{19}
Dém. par récurrence, c'est trivial si $\operatorname{card} I=0,1$ Le cas $\operatorname{card} I=2$ a été vu. Cas général \unsure{avec} $\operatorname{card} I\geq 3$, on choisit $J$ comme dessus. Il suffit de voir que les ens.
\[
\unsure{U|\deleted{=}U}_{J,i}\cap\Sigma_i^+,\ U_{J,i}\cap\Sigma_i^-
\]
ont la forme précédente, il suffit de le voir pour
\(\underset{\hspace{1.5em}\unsure{\vcenter{\hbox{%
% M01: drawing origin (1124,790) on the full scan; fixed local rectangle.
% The dark point and pale short lower-left trace have undetermined role/hand/origin.
\begin{tikzpicture}[baseline=(0,0),x=0.01cm,y=-0.01cm]
\path[use as bounding box] (0,0) rectangle (42,23);
\fill (35,7) circle[radius=0.02cm];
\draw[gray,line width=0.2pt] (6,18)--(6,20);
\end{tikzpicture}%
}}}}{\text{\unsure{l'une|l'un}}}\)
des deux.

Par hyp. de récurrence, \unsure{on a}
\[
U_{J,i}=X\setminus\mathop{\unsure{\coprod|\bigsqcup}}\limits_{\alpha\in \overset{\unsure{\cdot}}{\deleted{\unsure{J|\jmath}}}K_0}\widehat{\Sigma}_\alpha^+
\]
% M02: the point over the predecessor has full-scan centre about (989,1047).
% Its attachment is unresolved. The dotless j alternative prevents a second dot.
où
\[
\vcenter{\hbox{\unsure{%
% M03: pale trace left of brace; drawing origin (488,1276), pixel units.
% Its hand, character, role and surface origin remain undetermined.
\begin{tikzpicture}[baseline=(0,0),x=0.01cm,y=-0.01cm]
\path[use as bounding box] (0,0) rectangle (28,40);
\draw[gray!55,line width=0.18pt] (11,3)--(15,4);
\draw[gray!60,line width=0.23pt] (15,6).. controls (23,11) and (20,18)..(22,22)
  .. controls (22,28) and (13,32)..(8,36)--(1,36);
\end{tikzpicture}%
}}}\quad
\left\{\begin{aligned}
&\deleted{\Sigma_{\unsure{\alpha|0}}}K_0\subset J\\
&\Sigma_\alpha^+\in\pi_0(X\setminus E_\alpha)\qquad (\forall\,\alpha\in K_0)\\
&\widehat{\Sigma}_\alpha^+=\Sigma_\alpha^+\cup E_\alpha.
\end{aligned}\right.
\]
\unsure{On a|\deleted{On}On a} donc
\[
U=U_{J,i}\cap\Sigma_i^-=\unsure{\cdot}\,X\setminus\left(\coprod_{\alpha\in K_0}\widehat{\Sigma}_\alpha^+\ \cup\ \widehat{\Sigma}_i^+\right)
\qquad\underset{\unsure{\vcenter{\hbox{%
% M04: paired dark points; origin (1618,1684) on the full scan.
% These are copied as shapes, without an asserted punctuation or algebraic role.
\begin{tikzpicture}[baseline=(0,0),x=0.01cm,y=-0.01cm]
\path[use as bounding box] (0,0) rectangle (16,19);
\fill (10,6) ellipse[x radius=0.016cm,y radius=0.011cm];
\fill (5,13) ellipse[x radius=0.015cm,y radius=0.011cm];
\end{tikzpicture}%
}}}}{\vphantom{\Sigma}}
\]
% M05: the qualified point preceding X has full-scan centre about (728,1661).
% The second large join retains the first join's qualified primary reading.
\deleted{Soit} Au \unsure{deuxième|deuxieme|dernier} $=$, dans la \deleted{\unsure{somme|réunion}} réunion $\bigcup_{\alpha\in K_0}\widehat{\Sigma}_\alpha^+$, il suffit de se borner \unsure{à|:} sommer sur les \deleted{$\Sigma$} $\alpha\in K_0$ tels que $\widehat{\Sigma}_\alpha^+\not\subset\widehat{\Sigma}_i^+$ i.e. $\Sigma_\alpha^+\not\subset\Sigma_i^+$ soit $K_1$ cet ens., donc
\[
U=X\setminus\left(\bigcup_{\alpha\in K_1}\widehat{\Sigma}_\alpha^+\ \cup\ \widehat{\Sigma}_i^+\right)
\]
% The last lower limit touches the physical lower edge; no continuation is supplied.
% archive sheet number: 18
